\author{OTS}
\documentclass[12pt]{../../inc/mybibreport}

\setincpath{../../inc/}

\usepackage{bible_style}
\usepackage{bibeltext}
\usepackage{markierungen}

\graphicspath{{../../assets/images/}}

% Absatzformatierung
\setlength{\parindent}{0pt}
\setlength{\parskip}{0.6em}
\renewcommand{\arraystretch}{1.8}

\title{2. Petrus}
\author{Lothar Schmid}
\date{}

\begin{document}
\tableofcontents

\setdefault{SCHL}
\setversions{SCHL,ELB}

\chapter{Allgemeines}
\section{Verfassungszeit}
\section{Verfassungsort}
\section{Verfasser}
\section{Grund}

\chapter{Auslegung}
\section{Kapitel 1}
\subsection{Vers \textbf{1}}
\begin{spacing}{1.5}
\spa{0} \person[b, p]{Simon Petrus}, Knecht und Apostel \person[b, p]{Jesu Christi}, \\
    \spa{1}an die, \\
        \spa{2}\konj[b, p]{welche} den gleichen kostbaren Glauben wie wir \verbN[b, p]{empfangen haben}\\
        \spa{2}an die Gerechtigkeit unseres \person[b, p]{Gottes} \konj[b, p]{und} Retters \person[b, p]{Jesus Christus}:
\end{spacing}
\subsubsection{Auslegung}
    \begin{itshape}
        Der Name Simon ist der richtige Name von Petrus und Petrus ist der Name, der ihm Jesus gegeben hat.
    \end{itshape}
\subsubsection{Parallelstellen}
    \begin{itemize}
        \item \info[b, n]{Knecht} \bibleverse{Rom}(1:1); \bibleverse{Jak}(1:1)
        \item \info[b, n]{Glaube} \bibleverse{Eph}(4:5); \bibleverse{Tit}(1:4); \bibleverse{Jud}(1:3)
    \end{itemize}

\subsection{Vers \textbf{2}}
\begin{spacing}{1.5}
\spa{0}Gnade und Friede \\
    \spa{1} \verbN[b, p]{werde} euch mehr \konj[b, p]{und} mehr zuteil \\
    \spa{1}in der Erkenntnis \person[b, p]{Gottes} unseres Herrn \person[b, p]{Jesus}!
\end{spacing}
\subsubsection{Auslegung}
    \begin{itshape}
        Der Name Simon ist der richtige Name von Petrus und Petrus ist der Name, der ihm Jesus gegeben hat.
    \end{itshape}
\subsubsection{Parallelstellen}
    \begin{itemize}
        \item [Knecht] \bibleverse{Rom}(1:1); \bibleverse{Jak}(1:1)
        \item [Glaube] \bibleverse{Eph}(4:5); \bibleverse{Tit}(1:4); \bibleverse{Jud}(1:3)
    \end{itemize}

\subsection{Vers \textbf{3}}
\begin{spacing}{1.5}
Da seine göttliche Kraft uns alles geschenkt hat, was zum Leben und [zum Wandeln in] Gottesfurcht dient, durch die Erkenntnis dessen, der uns berufen hat durch [seine] Herrlichkeit und Tugend,
\end{spacing}
\subsubsection{Auslegung}
    \begin{itshape}
        Der Name Simon ist der richtige Name von Petrus und Petrus ist der Name, der ihm Jesus gegeben hat.
    \end{itshape}
\subsubsection{Parallelstellen}
    \begin{itemize}
        \item [Knecht] \bibleverse{Rom}(1:1); \bibleverse{Jak}(1:1)
        \item [Glaube] \bibleverse{Eph}(4:5); \bibleverse{Tit}(1:4); \bibleverse{Jud}(1:3)
    \end{itemize}






\end{document}